\clearpage\SourceLeaf{159}%
\thispagestyle{empty}
\pdfbookmark[0]{Reviews of His Own Writings}{notices-title}
\begin{center}
\vspace*{1in}{\Large REVIEWS\par}
\vspace{.4in}{\large OF HIS OWN\par}
\vspace{.4in}{\Large WRITINGS.\par}
\end{center}

\clearpage\SourceLeaf{161}%
\renewcommand{\BookRunningTitle}{Reviews of His Own Writings}
\pdfbookmark[1]{Theorematis arithmetici demonstratio nova. Review 1808 May 12}{notice-1808-05-12}
\begin{center}
\rule{.86\linewidth}{.3pt}\par
Göttingische gelehrte Anzeigen.\quad 1808 May 12.\par
\rule{.86\linewidth}{.3pt}
\end{center}

A paper submitted by Professor \textsc{Gauss} to the Royal Society of Sciences on January 15 of this year,
\begin{center}\emph{Theorematis arithmetici demonstratio nova},\end{center}
an account of whose contents we still owe here, concerns the famous fundamental theorem in the theory of quadratic residues, which plays so important a role both throughout \emph{higher arithmetic} and in the neighbouring parts of analysis. As is well known, an integer $a$ is called a \emph{quadratic residue} of the integer $b$ if there are numbers of the form $xx-a$ divisible by $b$, while in the opposite case $a$ is called a \emph{quadratic nonresidue} of $b$. The number $a$ may be positive or negative, whereas $b$ is always regarded as positive. Higher arithmetic teaches that all prime numbers $b$ for which a given number $a$ is a quadratic residue are included in certain linear forms, while other linear forms contain all the primes of which $a$ is a nonresidue. Thus, for example, $-1$ is a quadratic residue of all primes of the form $4n+1$ and a quadratic nonresidue of all primes of the form $4n+3$; further, $+2$ is a quadratic residue of all primes of the forms $8n+1$, $8n+7$, but a quadratic nonresidue of all primes of the forms $8n+3$, $8n+5$. There are infinitely many similar special theorems, but all can be derived by combining the two just mentioned
\SourceLeaf{162}%
with the following general theorem: Two distinct positive (odd) primes, $p$, $q$, always bear the \emph{same} relation to each other (that is, one is a quadratic residue or nonresidue of the other according as the other is a residue or nonresidue of the first) if either both have the form $4n+1$, or at least one does. Their mutual relation is opposite, however (that is, one is a nonresidue of the other if the latter is a residue of the former, and conversely), whenever \emph{both} have the form $4n+3$. This is the fundamental theorem mentioned above, which can be expressed in more than one form. The form chosen here is the one in which it is newly proved in Professor \textsc{Gauss}'s paper.

The most beautiful theorems of higher arithmetic, especially those discussed here, have the peculiarity that they are easily discovered by examining examples, whereas their proofs lie extremely well hidden and can be found only through very deep investigations. This is precisely what gives higher arithmetic the magical appeal that has made it the favourite field of the foremost mathematicians, quite apart from its inexhaustible wealth, in which it so far surpasses every other part of pure mathematics. The two special theorems mentioned above were already known to \textsc{Fermat}, who claimed also to possess their proofs. Whether he was mistaken in this we cannot decide, since he never published anything about them; but we can certainly regard it as possible, since there are several examples of self-deception among other great mathematicians, notably \textsc{Euler}, \textsc{Legendre}, and also \textsc{Fermat} himself. \textsc{Euler} gave the first proof of the first of these theorems. But this great mathematician did not succeed in proving the other, despite his devoted efforts over many years; it was left to \textsc{Lagrange} to fill that gap. Both mathematicians also proved various other special theorems, but a greater number, which they found by examining examples, always eluded their attempts at proof. It is, however, a remarkable twist of chance that neither mathematician discovered in this way the general fundamental theorem, which admits of so simple a statement. It was first presented, though in a somewhat different form, by \textsc{Legendre} in the \emph{Histoire de l'Académie des Sciences de Paris} 1785. Both there and later in his \emph{Essai d'une théorie des nombres},
\SourceLeaf{163}%
this excellent analyst sought to base the proof on very ingenious investigations, which nevertheless failed to reach the desired goal; and, unless we are mistaken, that goal could not have been reached by this route either.

The author of the paper to which this review is devoted entered the field of higher arithmetic at a time when all earlier work by other mathematicians in this subject was entirely unknown to him. It is chiefly to this circumstance that we must attribute his having followed a wholly distinctive course throughout. He found that fundamental theorem very early by examining examples, but only a whole year later, after many difficulties and unsuccessful attempts, did he succeed in finding the first completely rigorous proof, developed in the fourth section of his \emph{Disquisitiones Arithmeticae}. This proof, however, rests on very laborious and lengthy arguments. He subsequently found three other proofs, which are free of that inconvenience but instead require other very deep investigations, wholly different in content. One of these proofs is also given in the cited work, Art.~262; the other two will be made known in due course. The wish therefore remained that it might be possible to discover a shorter proof, independent of investigations of an unrelated kind. The author thus hopes that friends of higher arithmetic will take pleasure in a fifth proof, presented in the present paper in fewer than five pages, which seems to leave nothing to be desired in any respect. Given the concentrated brevity of this proof, we can admittedly give only an imperfect idea of its course here. Anything more would be all the more unnecessary, since Volume XVI of the \emph{Commentationes}, in which it has already been printed, will appear shortly.

\Needspace{30\baselineskip}
The proof rests on the following new theorem: Let $p$ denote a (positive odd) prime and $k$ any integer not divisible by $p$. Suppose further that, among the remainders obtained by dividing the $\frac12(p-1)$ products $k$, $2k$, $3k\dotsfour\frac12(p-1)k$ by $p$, there are in all $\mu$ remainders greater than $\frac12p$ (and hence $\frac12(p-1)-\mu$ less than $\frac12p$). Then $k$ will be a quadratic residue of $p$ if $\mu$ is even, but a quadratic nonresidue if $\mu$ is odd. The number $\mu$, which depends only on $k$
\SourceLeaf{164}%
and $p$, may be represented by the symbol $(k,p)$. A sequence of deductions that cannot be abridged now shows that, if $k$ and $p$ are two relatively prime odd numbers, $(k,p)+(p,k)+\frac14(k-1)(p-1)$ is always an \emph{even} number. It follows that, whenever $k$ and $p$ both have the form $4n+3$, one of the numbers $(k,p)$, $(p,k)$ must be even and the other odd; in all other cases, that is, whenever both numbers $k$ and $p$, or at least one of them, have the form $4n+1$, either .$(k,p)$, $(p,k)$ must both be even or both must be odd. Combined with the theorem above, this immediately gives the truth of the fundamental theorem.---The same route by which these results are obtained also gives, in the paper, a new proof of the two special theorems mentioned above. It is easy to show that $(-1,p)=\frac12(p-1)$, hence even or odd according as $p$ has the form $4n+1$ or $4n+3$; similarly, $(2,p)=\frac14(p-1)$ if $p$ has the form $4n+1$, and $(2,p)=\frac14(p+1)$ if $p$ has the form $4n+3$. Thus $(2,p)$ is even whenever $p$ has the form $8n+1$ or $8n+7$, but odd whenever $p$ has the form $8n+3$ or $8n+5$.
\BookRule{.28\linewidth}
